\chapter{Майорановское соответствие высшей степени на пакете \texorpdfstring{$H_{15}$}{H15}}

Статусная заморозка отделила текущий родитель $H_{15}/M_{35}$ от старой
вихревой архитектуры $H_{16}/M_{36}$. Теперь проверяется минимальная
совместимая альтернатива: может ли оператор высшей степени на одном левом
нейтринном канале получить единственный семейный тензор из уже имеющихся
голономии, проектора и следа.

\section{Граница известного механизма}

Без правого нейтрино первая калибровочно-инвариантная майорановская масса
имеет вид оператора Вайнберга
\begin{equation}
 \mathcal O_5=
 \frac{c_{ij}}{\Lambda}
 (\overline{L_i^c}\,\widetilde H^*)
 (\widetilde H^\dagger L_j).
 \label{eq:v5-h15-weinberg-operator}
\end{equation}
Слабое $SU(2)$-свёртывание единственно, но семейный коэффициент
$C=(c_{ij})$ является комплексной симметричной матрицей. Этот оператор был
введён С.~Вайнбергом как взаимодействие размерности пять, нарушающее
лептонное число на две единицы; см. DOI
\path{10.1103/PhysRevLett.43.1566}.

Таким образом, правильный тип оператора существует без добавления $N^c$,
но его семейная часть до наложения симметрии имеет шесть комплексных
коэффициентов. Следующий вопрос относится не к допустимости
\eqref{eq:v5-h15-weinberg-operator}, а к её сокращению текущей геометрией.

\section{Классификация при трёхциклической голономии}

Пусть трёхцикл на семейном триплете представлен матрицей
\begin{equation}
 C_3=
 \begin{pmatrix}
 0&0&1\\
 1&0&0\\
 0&1&0
 \end{pmatrix},
 \qquad C_3^3=I.
\end{equation}
Майорановская симметрия и ковариантность требуют
\begin{equation}
 C^T=C,
 \qquad
 C_3^T C C_3=C.
 \label{eq:v5-c3-majorana-invariance}
\end{equation}
Решение линейной системы \eqref{eq:v5-c3-majorana-invariance} двумерно.
Удобный базис задают спектральные проекторы
\begin{equation}
 P_0=\frac13(I+C_3+C_3^2)
 =\frac13
 \begin{pmatrix}
 1&1&1\\1&1&1\\1&1&1
 \end{pmatrix},
 \qquad Q=I-P_0.
 \label{eq:v5-c3-zero-projector}
\end{equation}
Поэтому общий инвариантный тензор имеет вид
\begin{equation}
 C=xP_0+yQ.
 \label{eq:v5-c3-majorana-family-tensor}
\end{equation}
На инвариантной линии его собственное значение равно $x$, а на
сопряжённой поперечной плоскости собственное значение $y$ имеет кратность
два.

Именно эта двумерность является известной границей циклических семейных
симметрий: циклическая перестановка ограничивает матрицу до циркулянтного
типа, но сама не задаёт все массы и смешивания; см.
\path{arXiv:hep-ph/0305243}.

Полная тетраэдрическая симметрия $A_4$ уменьшила бы коммутант до скаляров,
то есть оставила только $C\sim I$. Это устраняет относительный коэффициент,
но делает три семейных направления вырожденными и уничтожает роль
выбранной голономной оси.

\begin{nogocriterion}[Недостаточность одной $C_3$-ковариантности]
Без дополнительного спектрального ограничения трёхциклическая симметрия
оставляет два коэффициента $x,y$. Нормированный след
\begin{equation}
 \tau_3(C^*C)=\frac{|x|^2+2|y|^2}{3}
\end{equation}
задаёт только длину и не фиксирует их отношение.
\end{nogocriterion}

\section{Каноническое сжатие на нулевую ветвь}

Предыдущий голономный гейт дал ровно одну нулевую семейную ветвь. Её
проектор не выбирается вручную: формула \eqref{eq:v5-c3-zero-projector}
является полиномом от самой голономии. Поэтому низкоэнергетическое сжатие
майорановского соответствия имеет вид
\begin{equation}
 P_0 C P_0=xP_0.
 \label{eq:v5-zero-compressed-majorana}
\end{equation}
Пространство семейных направлений после сжатия одномерно. Тем самым
исчезает прежняя свобода выбора вектора $(1,1,1)$: он восстановлен как
единственная собственная линия с нулевым голономным импульсом.

Общий след $M_{35}$ также фиксирует относительную кратность этого канала.
Вес наблюдаемого угла равен $3/7$, а нормированный ранг $P_0$ внутри
триплета равен $1/3$. Следовательно,
\begin{equation}
 w_0=\frac37\cdot\frac13=\frac17,
 \qquad
 w_\perp=\frac37\cdot\frac23=\frac27.
 \label{eq:v5-h15-zero-line-weight}
\end{equation}
Эти числа являются весами чтения, а не величиной нейтринной массы.

\begin{proposition}[Одномерный семейный проход]
Если майорановское соответствие высшей степени сначала ограничивается на
уже выведенную голономную нулевую ветвь, его семейное направление
единственно и равно $P_0$. Никакой дополнительный семейный вектор или
относительный коэффициент не требуется. Общая амплитуда $x/\Lambda$ при
этом остаётся невыведенной.
\end{proposition}

\section{Подсказка фермионного спектрального действия}

В первичной литературе существует конструкция, непосредственно близкая к
этому разрыву. М.~Сакеллариаду и А.~Ситарз показали, что в
пятнадцатимерной почти коммутативной геометрии без правого нейтрино член,
квадратичный по полю Хиггса и майорановский по левым нейтрино, может
возникать в следующем порядке фермионного спектрального действия при
использовании нескалярной спектральной функции; см.
\path{arXiv:1903.09149}.

Это важный положительный ориентир: требуемая высшая степень совместима с
полевым составом $H_{15}$. Но он ещё не замыкает текущую программу по трём
причинам:
\begin{enumerate}
 \item нескалярный эндоморфизм спектральной функции является дополнительным
 объектом и не выведен из родителя $M_{35}$;
 \item коэффициент зависит от момента функции отсечения и масштаба
 $\Lambda$, которые $C_3$ и след не фиксируют;
 \item эффективный массовый член не выводит радиальный дефектный профиль и
 длину локализации.
\end{enumerate}

Следовательно, литературная подсказка снимает запрет существования
оператора, но переводит вопрос к общей фермионной мере.

\begin{theorem}[Частичный проход соответствия на $H_{15}$]
На текущем пакете $H_{15}$ существует допустимый калибровочный тип
майорановского оператора высшей степени. Одна $C_3$-ковариантность оставляет
два семейных коэффициента, однако каноническое спектральное сжатие на
единственную нулевую ветвь фиксирует семейное направление до общей
амплитуды. Не выведены сама амплитуда, момент нескалярного фермионного
спектрального действия и объёмная локализация.
\end{theorem}

\section{Следующий критерий}

Следующий гейт должен проверить уже не семейный вектор, а коэффициент
\eqref{eq:v5-zero-compressed-majorana}: можно ли получить нескалярный
фермионный спектральный эндоморфизм, его момент и масштаб из той же меры,
которая дала веса $M_{35}$, без подстановки нейтринной массы. Если остаётся
произвольная функция отсечения или независимый масштаб, ветвь закрывается
как условный эффективный оператор.

\begin{protocol}
\begin{enumerate}
 \item калибровочный тип оператора Вайнберга на $H_{15}$:
 \textbf{допустим};
 \item число слабых свёртываний: \textbf{одно};
 \item размерность симметричных $C_3$-инвариантных семейных тензоров:
 \textbf{два};
 \item проектор единственной нулевой ветви: \textbf{каноничен};
 \item семейное направление после сжатия: \textbf{одномерно};
 \item вес нулевой линии в следе $M_{35}$: \textbf{$1/7$};
 \item общая майорановская амплитуда: \textbf{не выведена};
 \item дефектный профиль: \textbf{не выведен};
 \item физическое замыкание: \textbf{не достигнуто}.
\end{enumerate}
\end{protocol}

Машинный протокол сохранён в
\path{s2t_v5_h15_majorana_pairing_correspondence_gate_results.json}.